%% This document can only be compiled correctly with LuaLaTeX!
%% Go to Menu > Settings > Compiler and select LuaLaTex from the dropdown menu. Then hit recompile to embed the marked expressions into metadata.
\documentclass{article}
%% For more information visit https://github.com/Christof93/SciKGTeX 
%% © Christof Bless, 2024

\usepackage[colorlinks=true, allcolors=blue]{hyperref}
\usepackage[compatibility]{scikgtex}
% add title to metadata
\metatitle*{Further Research into this Exciting Topic}
% add author to metadata
\metaauthor*{John D. Doe}
\metaauthor*{Helen R. Keller}
\metaauthor*{Ahmet Saladin}

%% This information markedd with the starred variant of the SciKGTeX commands is only visible in the document metadata and not in the text.
%% a URI is used to link the topic with a corresponding uniquely identifiable pendant in the web.
\researchfield*{\uri{https://www.orkg.org/orkg/resource/R12259}{antibiotic therapy}}
\contribution*{Uses drug}{Soandso}

\title{
  Further Research into this Exciting Topic
}
\author{
  John D. Doe
  \and
  Helen R. Keller
  \and
  Ahmet Saladin
}

\begin{document}
\maketitle

\begin{abstract}
%% In the abstract we denote research problem, objective, method result and conclusion of the study.
%% Notice that these commands can be used anywhere in the document and also more than once.
%% The information inside the brackets ends up in the metadata in the end.
The role of \researchproblem{antibiotic therapy in managing acute bacterial sinusitis (ABS) in children} is controversial.
The purpose of this study was to \objective{determine the effectiveness of the drug called Soandso}.
This was a \method{randomized, double-blind, placebo-controlled study}.
\result{Children receiving the antibiotic were more likely to be cured (50\% vs 14\%).}
ABS is a common complication of viral upper respiratory infections. \conclusion{Soandso results in significantly more cures and fewer failures than placebo.}
\end{abstract}

\section{Introduction}
Here we introduce our research..
\section{Methods}
%% additional to the mandatory commands found in the abstract text, we can also define custom properties of the contribution. For example here we add the sample size of the study. Note how only the second part of the command is rendered.
For our study we use a sample size of \contribution{sample size}{1500} participants.
\end{document}
